% Complete source: originales/cuaderno/REORIENTACION_GENEALOGICA_Y_RESULTADOS.md,
% §§2--7. Management and provenance records are preserved separately.
\Needspace{16\baselineskip}
\section{Genealogical regularity, operator reconstruction, and the geometry of action}
\label{md:restop:seccion}

\subsection{Production, compatibility, and reconstruction}
\label{md:restop:jerarquia}

The primary phenomenon is a dynamics that preserves provenance while producing
new observations. Three operations with distinct mathematical roles
are brought together:
\begin{enumerate}
 \item \emph{Production.} The APP--TRIT--TPK update prolongs the state
 and its correlated regions. The phase may return while memory
 advances.
 \item \emph{Compatibility.} The regional, angular, action,
 and incidence readouts obey common relations. Their variation must respect
 those relations in order to correspond to the same state.
 \item \emph{Reconstruction.} Retaining complementary information makes it possible
 to recover earlier values, products of operators, and differences
 between orders of composition.
\end{enumerate}
A comparison between procedures for obtaining $\pi$ must attend to these
three levels. The specific HMT thesis lies in the joint construction
and its relations, in addition to the production of numerical expansions.
A decimal coincidence or the periodicity of a phase does not, by itself,
describe the scope of the construction. The distinction identifies the
mathematical object relevant to a historical comparison; the focal development
does not, by itself, establish worldwide priority.

\subsection{Vacancy regimes and return of the capacity phase}
\label{md:restop:vacancias}

\subsubsection{Capacity calendar and tritic selector}
\label{md:restop:antecedentes-capacidad}

The capacity readout receives the internal compactification $729/1000$ and
distinguishes the refinement index from the capacity index. With the
convention of Section~\ref{vac-capacidad-trit},
\begin{equation}
 \delta=\log_{10}(10/9),\qquad
 p_j=\left\lfloor\frac{j}{\delta}\right\rfloor,\qquad
 v_j=p_j-j,\qquad j\geq1.
 \label{md:restop:capacidad}
\end{equation}
The intervals between events are
\begin{equation}
 h_j=p_{j+1}-p_j\in\{21,22\},\qquad
 s_j=22-h_j.
 \label{md:restop:intervalos}
\end{equation}
The tritic class obeys
\begin{equation}
 [v_{j+1}]_3=[v_j-s_j]_3.
 \label{md:restop:clase-tritica}
\end{equation}
A short interval changes the class; a long one preserves it.
The $6/7$ spacings describe the plateaux of this induced regime.
The construction retains its type as an observation factor of the TPK state;
it neither replaces that state nor identifies those spacings with charge
signs or with digits of $\pi$ or $e$.
The canonical selector is
$M_{\mathrm{ph}}(1+[v_j]_9)$, whose pattern depends on the class modulo three.

\subsubsection{Restoration of the local regime and displacement modulo nine}
\label{md:restop:consecuencia-capacidad}

Let
\begin{equation}
 \beta_{\mathrm{cap}}=22-\delta^{-1},\qquad
 q_k=\left\lfloor\frac{k}{\beta_{\mathrm{cap}}}\right\rfloor,\qquad k\geq1.
 \label{md:restop:pendiente-inducida}
\end{equation}
The positions $q_k$ are the indices of the short intervals. Indeed,
\begin{equation}
 p_j=22j-\lceil j\beta_{\mathrm{cap}}\rceil,\qquad
 s_j=\lceil(j+1)\beta_{\mathrm{cap}}\rceil
     -\lceil j\beta_{\mathrm{cap}}\rceil.
 \label{md:restop:eventos-cortos}
\end{equation}
Irrationality of $\delta$ follows from the valuations at $2,3,5$
in the equality that $(10/9)^q=10^p$ would impose; it implies irrationality of
$\beta_{\mathrm{cap}}$. Consequently, positive indices have no
endpoint ambiguity. The discontinuity corresponding to the integer $k$
occurs at $j=\lfloor k/\beta_{\mathrm{cap}}\rfloor$.

\begin{proposition}
\label{md:restop:proposicion-regimen}
Between the events indexed by $q_k$ and $q_{k+3}$, taking the intervals with
indices in $[q_k,q_{k+3})$, the TRIT class returns while the capacity phase
modulo nine changes.
\end{proposition}
\begin{proof}
There are exactly three short intervals in this segment.
The bounds $1/7<\beta_{\mathrm{cap}}<3/20$ give
$20<3/\beta_{\mathrm{cap}}<21$, hence
\begin{equation}
 m_k=q_{k+3}-q_k\in\{20,21\}.
 \label{md:restop:duracion-tritica}
\end{equation}
The sum of lengths and the capacity gain are, respectively,
\begin{equation}
 \Delta p=22m_k-3\in\{437,459\},\qquad
 \Delta v=21m_k-3\in\{417,438\}.
 \label{md:restop:ganancias}
\end{equation}
In both cases $\Delta v\equiv0\pmod3$. By contrast,
\begin{equation}
 \boxed{\Delta v\equiv3\ \text{or}\ 6\pmod9.}
 \label{md:restop:desplazamiento-nonadico}
\end{equation}
This proves the statement.
The bounds on $\beta_{\mathrm{cap}}$ have integer certificates:
\begin{equation}
 \begin{aligned}
 10^{146}>9^{153}&\ \Longrightarrow\
 \delta>\frac7{153},\\
 10^{417}<9^{437}&\ \Longrightarrow\
 \delta<\frac{20}{437}.
 \end{aligned}
 \label{md:restop:certificados-cotas}
\end{equation}
Monotonicity of $22-1/x$ yields
$1/7<\beta_{\mathrm{cap}}<3/20$. Thus the bounds rest on
integer inequalities, not on substitution of a rounded decimal.
\end{proof}

The selector value
$\tau_j=M_{\mathrm{ph}}(1+[v_j]_9)$ returns because it depends on the class
modulo three. This return remains distinct from the full TRIT
state and from a chronological turn of $\Gamma_9$.
Both durations occur infinitely often by density of the irrational rotation
with slope $1/\beta_{\mathrm{cap}}$.

Regularity has distinct levels: the local regime is restored
while the capacity phase remains displaced. The nonadic return of the
state also has its own memory increment. Each object
therefore retains its own notion of return and its own duration.
The numbers $437$ and $459$ are durations in this induced observation,
without being assigned the status of additional universal constants.
The result explicitly assembles the preceding formulae; presenting it
as a focal consequence does not assert historical priority
over developments that have not been exhaustively compared.

\subsection{The contribution of memory to operator composition}
\label{md:restop:memoria}

\subsubsection{Enlarged representation of readout and memory}
\label{md:restop:estructura-partida}

The readout--memory representation is obtained by separating eight
components from a ninth. In a single fibre, with identity reference
transport, the internal transport $C$ induces
\begin{equation}
 T_C=\frac{8I+C}{9},\qquad
 D_C=\frac{\sqrt8(C-I)}{9},\qquad
 R_C=\frac{I+8C}{9},
 \label{md:restop:bloques}
\end{equation}
\begin{equation}
 \mathcal U_C=
 \begin{pmatrix}
 T_C&D_C\\D_C&R_C
 \end{pmatrix}.
 \label{md:restop:transporte-ampliado}
\end{equation}
The coupling arises from the change of basis in the nonadic decomposition.
For transports preserving the declared form, it preserves the enlarged
form. The composition identity is
\begin{equation}
 \mathcal U_{C_2}\mathcal U_{C_1}
 =\mathcal U_{C_2C_1},\qquad
 T_2T_1+D_2D_1=\frac{8I+C_2C_1}{9},
 \label{md:restop:composicion-ampliada}
\end{equation}
with $T_i=T_{C_i}$ and $D_i=D_{C_i}$.
The second term corresponds to memory produced in the first
operation and reintroduced in the second. Setting it aside changes the
procedure: it produces $T_2T_1$ instead of the readout of the enlarged
composition.

\subsubsection{Partition of sensitivity to the order of composition}
\label{md:restop:conmutadores}

\begin{proposition}
\label{md:restop:proposicion-orden}
With the convention $[A,B]=AB-BA$,
\begin{equation}
 [T_2,T_1]=\frac1{81}[C_2,C_1],\qquad
 [D_2,D_1]=\frac8{81}[C_2,C_1].
 \label{md:restop:reparto-conmutadores}
\end{equation}
Consequently, the readout difference between the two orders, with
memory reintroduced, is
\begin{equation}
 \boxed{(T_2T_1+D_2D_1)-(T_1T_2+D_1D_2)
 =\frac19[C_2,C_1].}
 \label{md:restop:orden-completo}
\end{equation}
\end{proposition}
\begin{proof}
On expanding the products of $T$, the constant and linear terms
cancel in the subtraction, leaving the commutator divided by $81$.
Expanding the products of $D$ gives the same cancellation, with a
factor of $8/81$. Their sum is $1/9$.
\end{proof}

The re-entry mechanism belongs to the enlarged representation.
Comparing the two orders makes one consequence explicit:
in this representation, the memory contribution to the antisymmetric
response is eight times that retained by composing isolated
readouts. The full response is nine times the latter. These
ratios concern response operators, not percentages
of energy.

Reconstruction preserves noncommutativity and its readout intensity.
In this case, isolated compression attenuates a nonzero commutator without
annihilating it. An application to physical forces also retains the interaction
operators and the corresponding work functional; the algebraic
identity does not identify them by itself.

\subsubsection{Operator comparison and nonadic normalization}
\label{md:restop:contraste}

For the general algebraic pattern $(n-1)+1$, the same calculations give the
factors $1/n^2$, $(n-1)/n^2$, and $1/n$. The nonadic architecture fixes
$n=9$. The composition law and the particular value of
its normalization are thus distinguished.

When the effective operations of an application are represented by
$C_1,C_2$, procedures that set aside or reintroduce memory
predict different order responses. The comparison concerns operational
composition and the accessibility of memory, preserving
the coefficients of the representation.

\subsubsection{Readout and memory norms under central inversion}
\label{md:restop:treinta-y-dos}

For $C=-I$, one obtains
\begin{equation}
 T^*T=\frac{49}{81}I,\qquad
 D^*D=\frac{32}{81}I.
 \label{md:restop:normas}
\end{equation}
The numerator $32$ comes from $8|{-1}-1|^2$.
In the nine memory components, the norm includes
$8^2+8=72$, with denominator $27^2$. The square $64$ enters
that calculation together with the other eight contributions.
Identification with a region of $64$ cells requires the region's
incidence map and is not obtained by doubling the numerator $32$.
A geometric readout of an $8\times8$ support against a
$9\times9$ support therefore retains the problem of its connection to the
particular application; an isolated numerical equality does not replace it.

\subsection{Action character and complete reciprocity coordinate}
\label{md:restop:h5}

\subsubsection{Genealogical coefficients and regional continuation}
\label{md:restop:coeficientes}

The central matrix and its invariants are
\begin{equation}
 B_c=\begin{pmatrix}7&2\\2&7\end{pmatrix},\qquad
 \lambda_+=9,\qquad\lambda_-=5.
 \label{md:restop:matriz-central}
\end{equation}
The length-four orbit of the step-three shift, the
length-twelve calendar, and the census quotient $104976/1944=54$ also enter.
The coefficient relations are
\begin{equation}
 \begin{aligned}
 \frac9{16}&=\frac{\lambda_+}{4^2},&
 \frac59&=\frac{\lambda_-}{\lambda_+},\\
 \frac7{48}&=\frac{\operatorname{tr}B_c/2}{4\cdot12},&
 \frac1{54}&=\frac{1944}{104976}.
 \end{aligned}
 \label{md:restop:coeficientes-h5}
\end{equation}
The assignment of degrees and signs belongs to the definition of the character;
its isolated cardinalities do not determine it. The complete readout is
\begin{equation}
 H_5=\frac12\sqrt{\frac{1000\alpha}{\varphi}}-\alpha
 +\frac9{16}\alpha^2-\frac59\alpha^3
 +\frac7{48}\alpha^4-\frac1{54}\alpha^5,
 \label{md:restop:h5-completo}
\end{equation}
\begin{equation}
 \begin{aligned}
 D_A&=\exp(-100\pi\alpha/9),\\
 \mathcal C_\pi&=169\alpha^6+
 \frac{D_A\alpha^7}{1-D_A\alpha},\\
 \eta_{\mathrm{ret}}&=H_5-\frac{90}{\pi}\mathcal C_\pi.
 \end{aligned}
 \label{md:restop:retorno-regional}
\end{equation}
The fraction preserves the infinite continuation defined by
\begin{equation}
 \mathscr R(z)=D_Az^7+D_Az\mathscr R(z),
 \label{md:restop:renovacion}
\end{equation}
with unit initial weight. The continuation remains intact in the
readout; a truncated polynomial does not replace that object. The construction
of the coefficients and the renewal are assembled in
Section~\ref{sec:revision-planck}.

\subsubsection{Classification of oriented radial reciprocity}
\label{md:restop:geometria}

The relations $A_{\deg}=1000\alpha$,
$C^*_{\deg}=2(\eta_{\mathrm{ret}}+\alpha)$ and the radial recovery map of
Theorem~\ref{rec:thm:accion} are retained. To keep both orientations, let
\begin{equation}
 \epsilon\in\{+1,-1\},\qquad
 q_{\mathrm{rad}}=R_\star^4>0,\qquad
 \Xi_{\mathrm{rad}}=
 \epsilon\frac{1-q_{\mathrm{rad}}}{1+q_{\mathrm{rad}}}.
 \label{md:restop:par-radial}
\end{equation}
The recovery map \eqref{rec:eq:eta-orientada} is written
\begin{equation}
 \frac{\eta_{\mathrm{ret}}}{\alpha}=500\Xi_{\mathrm{rad}}-1.
 \label{md:restop:recuperador}
\end{equation}
The involution
$(\epsilon,q_{\mathrm{rad}})
 \mapsto(-\epsilon,q_{\mathrm{rad}}^{-1})$
preserves that action. The coordinate assembling this structure is
\begin{equation}
 \boxed{\chi_{\mathrm{rad}}=\epsilon\log q_{\mathrm{rad}}.}
 \label{md:restop:chi}
\end{equation}

\begin{proposition}
\label{md:restop:proposicion-orbitas}
$\chi_{\mathrm{rad}}$ classifies exactly the orbits of this involution
on $\{\pm1\}\times\mathbb R_{>0}$.
\end{proposition}
\begin{proof}
Simultaneously changing the sign and inverting $q_{\mathrm{rad}}$ preserves
$\chi_{\mathrm{rad}}$. If two pairs have the same value of
$\chi_{\mathrm{rad}}$, equal signs imply equal
$q_{\mathrm{rad}}$; with opposite signs, their fourth powers of the radii are
reciprocal. These are exactly the two cases of belonging to the same
orbit. The argument includes $q_{\mathrm{rad}}=1$, whose orbit retains
both orientations.
\end{proof}

The identity
$(1-q)/(1+q)=-\tanh(\log q/2)$ gives
\begin{equation}
 \Xi_{\mathrm{rad}}
 =-\tanh(\chi_{\mathrm{rad}}/2),\qquad
 \boxed{\eta_{\mathrm{ret}}
 =\alpha[-500\tanh(\chi_{\mathrm{rad}}/2)-1].}
 \label{md:restop:accion-chi}
\end{equation}
Compatibility with the full character takes the form
\begin{equation}
 \boxed{\chi_{\mathrm{rad}}
 =-2\operatorname{artanh}
 \left(\frac{H_5+\alpha-(90/\pi)\mathcal C_\pi}{500\alpha}\right).}
 \label{md:restop:chi-caracter}
\end{equation}
The term $-\alpha$ in $H_5$ cancels the $+\alpha$ from the
angular contrast before this readout. The denominator of the
regional continuation remains complete. The equality receives the
preceding HMT outputs and readout operators and assembles their two representations;
it does not define $\alpha$ retroactively.

On the branch of positive action and finite radius,
\begin{equation}
 0<\eta_{\mathrm{ret}}<499\alpha,\qquad
 \chi_{\mathrm{rad}}
 <-2\operatorname{artanh}(1/500).
 \label{md:restop:dominio-positivo}
\end{equation}
At fixed $\alpha$,
\begin{equation}
 \frac{\partial\eta_{\mathrm{ret}}}{\partial\chi_{\mathrm{rad}}}
 =-250\alpha\,\operatorname{sech}^2(\chi_{\mathrm{rad}}/2)<0.
 \label{md:restop:derivada-accion}
\end{equation}
Thus the dimensionless action coefficient uniquely determines the
oriented radial class in this representation, and conversely.
For joint variations,
\begin{equation}
 d\eta_{\mathrm{ret}}
 =\frac{\eta_{\mathrm{ret}}}{\alpha}\,d\alpha
 -250\alpha\,\operatorname{sech}^2(\chi_{\mathrm{rad}}/2)
 \,d\chi_{\mathrm{rad}}.
 \label{md:restop:diferencial-accion}
\end{equation}

The character $H_5$, the regional correction, and the oriented radial anisotropy
thus satisfy a compatibility relation rather than constituting three
independent parameters. The analytic character and the geometry of
reciprocity represent the same return section.
This structural level of the Planck equation remains explicit
before action is composed with a frequency.

The action decade retains its determinantal genealogy, independently
of this reparametrization:
\begin{equation}
 \hbar_{\mathrm{ret}}
 =10^{-34}\mathcal U_S\eta_{\mathrm{ret}},\qquad
 \hbar_{\mathrm{pre}}=10^{-34}\mathcal U_SH_5.
 \label{md:restop:secciones-accion}
\end{equation}
The action per turn $h=2\pi\hbar$ and the thermal readout
\begin{equation}
 k_B\Theta_{\mathrm{clk}}\log3=\frac{h}{T_{108}}
 \label{md:restop:lectura-termica}
\end{equation}
are also retained within their scope. The physical clock of a realization retains its construction:
the coordinate $\chi_{\mathrm{rad}}$ classifies this reciprocity, not
the totality of TPK histories. Classification by
$\chi_{\mathrm{rad}}$ and the differential relation assemble the existing
representations of action; they retain their character as a focal formalization and do
not constitute a second independent derivation of Planck's constant.

\subsection{Phase circuits, memory, and conservativity}
\label{md:restop:conservatividad}

Nonadic holonomy restores the phase and increments memory, as described in
Section~\ref{mono-ampl:vueltas}. If a path closed in the phase projection
increases the integer counter, that increment cannot be written as the
difference of a potential depending only on phase.

The proof follows by telescoping: the sum
\begin{equation}
 \sum_j\bigl[F(g_{j+1})-F(g_j)\bigr]=0
 \label{md:restop:potencial-fase}
\end{equation}
around a phase circuit contrasts with the memory increment, equal
to one per turn. On the enlarged state, the counter $m$ records
that increment. The path closed in the projection therefore
remained open in the full state.

The study of conservativity begins by declaring the space in which
the path closes and the coordinates participating in the
dynamics. Noncommutativity and nonconservativity are different
properties. Section~\ref{md:restop:memoria} controls the order of
operations; this section distinguishes a projected circuit from a full
circuit. To study a particular force, the action or
work functional must act on the same state space and the same transport.

Memory growth does not by itself assign thermodynamic
production. Section~\ref{md:cron:crecimiento} proves that
block information can grow with zero entropy rate.
Operator reconstruction adds a clarification: preserving information
may mean preserving which transformations can be composed,
in addition to retaining the stored symbols.

\subsection{Dynamical compatibility between observations}
\label{md:restop:sintesis}

The unit of study is dynamical compatibility between observations,
with memory as part of the operation. The value of a constant
constitutes a readout; its structure contains constraints on other
readouts and on admissible transformations. Prolonging the
analysis consists in obtaining relations that prevent value, orientation,
memory, form, and action from being varied independently.

The three focal developments follow this same axis:
\begin{itemize}
 \item A complete cycle of TRIT regimes still retains a
 nonadic capacity displacement: local regularity does not exhaust
 return.
 \item Memory re-entry changes exactly the response to the order
 of two operations: the history participates causally in the composition.
 \item The character $H_5$ and the oriented radial geometry mutually
 determine their return section: action admits a recoverable
 structural representation.
\end{itemize}
These consequences continue the genealogy of $\pi,\varphi,e,\alpha$ and
make explicit what it means to preserve it when passing to operations and
action. The subsequent composition assembles these levels through the
common maps already constructed, maintaining the distinction between their indices and
memories even when they share a name.
